% from ega4-16-fr.tex
\oldpage[IV-4]{36}
suffit de voir que l'on a $d(xy)-x.dy-y.dx=0$, ce qui exprime que $d$ est une
dérivation.

Il reste à prouver que $d$ vérifie la condition (i). Or, le carré d'une
antidérivation est une dérivation (Bourbaki, \emph{loc. cit.}), et comme
$\Omega_{B/A}^{\bullet}$ est engendré par $\Omega_{B/A}^1$ comme
$B$-algèbre, il suffit de vérifier que $d(dz)=0$ pour $z\in B$ et
$z\in\Omega_{B/A}^1$; dans le premier cas, cela résulte de la formule
(16.6.2.3) avec $g=1$; dans le second, on peut se borner au cas où $z=g.df$
avec $f$, $g$ dans $B$, et alors, en vertu de (16.6.2.1) et (16.6.2.3), on a
\[
  d(d(g.df))=d(dg\wedge df)
  =(d(dg))\wedge(df)-(dg)\wedge(d(df))=0.
\]
\hfill C.Q.F.D.

\begin{definition}[16.6.3]
\EGAReaderTarget{ega4.c04.def.16-6-3}{IV.16.6.3}
L'antidérivation $d$ définie dans (16.6.2) (aussi notée $d_{X/S}$) est
appelée \emph{différentielle extérieure sur $X$} (relativement à $S$).
\end{definition}

\begin{proposition}[16.6.4]
\label{IV.16.6.4}
Pour tout changement de base $g:S'\to S$, si l'on pose
$X'=X\times_S S'$, l'homomorphisme canonique
\[
  \Omega_{X/S}^{\bullet}\otimes_S S'
  \longrightarrow\Omega_{X'/S'}^{\bullet}
  \tag{16.6.4.1}
  \label{IV.16.6.4.1}
\]
déduit de l'isomorphisme (16.5.9.1) est bijectif. En outre, si $s$ est une
section de $\Omega_{X/S}^{\bullet}$ au-dessus d'un ouvert $U$ de $X$,
$s\otimes1$ son image réciproque, section de $\Omega_{X'/S'}^{\bullet}$
au-dessus de l'image réciproque $U'$ de $U$ dans $X'$, on a
$d_{X'/S'}(s\otimes1)=d_{X/S}(s)\otimes1$.
\end{proposition}

La première assertion est immédiate, la formation de l'algèbre extérieure
d'un module permutant avec toute extension de l'anneau des scalaires. Pour
prouver la seconde, on peut, en vertu de (16.6.2.2), se borner au cas où
$s\in\Gamma(U,\mathcal{O}_X)$, et dans ce cas l'assertion a déjà été prouvée
(16.4.3.7).

\begin{env}[16.6.5]
\label{IV.16.6.5}
Supposons que $\Omega_{X/S}^1$ soit un $\mathcal{O}_X$-Module localement libre
de rang $n$ en un point $x$, de sorte qu'il y a $n$ sections
$s_i\in\Gamma(U,\mathcal{O}_X)$ telles que les $ds_i$ forment une base du
$\Gamma(U,\mathcal{O}_X)$-module $\Gamma(U,\Omega_{X/S}^1)$ (16.5.8).
Alors, pour tout entier $p\geq1$, les $p$-différentielles
$ds_{i_1}\wedge ds_{i_2}\wedge\cdots\wedge ds_{i_p}$ (pour
$i_1<i_2<\cdots<i_p$, éléments de $[1,n]$) forment une base de
$\binom{n}{p}$ éléments de $\Gamma(U,\Omega_{X/S}^p)$ sur
$\Gamma(U,\mathcal{O}_X)$. En outre la formule (16.6.2.2) montre que pour
toute section $g\in\Gamma(U,\mathcal{O}_X)$, on a
\[
  d(g.ds_{i_1}\wedge ds_{i_2}\wedge\cdots\wedge ds_{i_p})
  =\sum_k(-1)^r\frac{\partial g}{\partial s_k}
  ds_{i_1}\wedge\cdots\wedge ds_{i_r}\wedge ds_k\wedge
  ds_{i_{r+1}}\wedge\cdots\wedge ds_{i_p}
  \tag{16.6.5.1}
  \label{IV.16.6.5.1}
\]
où, au second membre, $k$ parcourt l'ensemble des $n-p$ indices distincts des
$i_h$, $i_r$ étant le plus grand indice $<k$.

On notera que la relation $d(dg)=0$ pour toute section
$g\in\Gamma(U,\mathcal{O}_X)$ s'exprime sous la forme
\[
  D_i(D_jg)=D_j(D_ig)\qquad\text{pour $i\neq j$};
\]
autrement dit, les dérivations $D_i$ définies dans (16.5.7) sont deux à deux
permutables.
\end{env}

\subsection{\texorpdfstring{Les $\mathscr{P}_{X/S}^n(\mathscr{F})$.}{Les P\string_\string{X/S\string}\string^n(F).}}
\EGAReaderTarget{ega4.c04.subsec.16-7}{IV.16.7}

\begin{env}[16.7.1]
\label{IV.16.7.1}
Soient $f:X\to S$ un morphisme de préschémas, $\mathscr{F}$ un
$\mathcal{O}_X$-Module. Désignons par $X_{\Delta_f}^{(n)}$ le \emph{$n$-ème
voisinage infinitésimal de $X$} pour le morphisme diagonal
